\chapter{Thực Nghiệm và Đánh Giá}
\label{chap:experiment}

\section{Môi Trường Thực Nghiệm}
\label{sec:env}

Các thí nghiệm được thực hiện trên nền tảng Kaggle Notebooks với cấu hình:
\begin{itemize}
  \item \textbf{CPU}: Intel Xeon (2 cores, 13GB RAM)
  \item \textbf{GPU}: NVIDIA Tesla T4 (16GB VRAM)
  \item \textbf{Framework}: PyTorch 2.6, scikit-learn 1.4, Python 3.12
  \item \textbf{Seed}: 42 (tái tạo được kết quả)
\end{itemize}

\section{Bộ Dữ Liệu Mô Phỏng}
\label{sec:dataset}

Do thiếu dữ liệu thực từ người dùng thật, luận văn sử dụng bộ dữ liệu mô phỏng được tạo ra bằng pipeline IRT-based. Quá trình sinh dữ liệu mô phỏng hành vi ứng viên theo ba bước: (1) tạo hồ sơ người dùng với vector năng lực theo phân phối chuẩn, (2) sinh chuỗi tương tác dựa trên mô hình IRT 3PL với noise $\sigma = 0.30$, và (3) sinh self-rating log làm quality label cho DKT/SAKT.

\subsection{Thống kê tổng quan}

\begin{table}[H]
  \centering
  \caption{Thống kê bộ dữ liệu mô phỏng}
  \label{tab:dataset}
  \begin{tabular}{lr}
    \toprule
    \textbf{Thuộc tính} & \textbf{Giá trị} \\
    \midrule
    Tổng số người dùng & 1{,}000 \\
    \quad -- Data Analyst (DA) & 334 \\
    \quad -- Data Scientist (DS) & 333 \\
    \quad -- Data Engineer (DE) & 333 \\
    Tổng số câu hỏi & 1{,}578 \\
    Số Knowledge Components (KC) & 17 \\
    Tổng số tương tác & 40{,}140 \\
    Trung bình tương tác/người & 40.1 \\
    \midrule
    \textbf{Phân chia dữ liệu} & \\
    \quad Train & 700 users / 28{,}164 tương tác \\
    \quad Validation & 150 users / 5{,}911 tương tác \\
    \quad Test & 150 users / 6{,}065 tương tác \\
    \midrule
    Tỷ lệ nhãn positive (Pass) & 82.5\% \\
    \bottomrule
  \end{tabular}
\end{table}

\subsection{Kiểm tra tính toàn vẹn dữ liệu}

Trước khi huấn luyện, bộ dữ liệu được kiểm tra qua 7 tiêu chí:
\begin{enumerate}
  \item \textbf{Data leakage}: không có user ID nào xuất hiện ở cả train và test.
  \item \textbf{Label distribution}: positive rate train (82.7\%) $\approx$ test (82.2\%), chênh lệch $< 1\%$.
  \item \textbf{Sequence length}: độ dài trung bình 40 bước, phân phối đồng đều.
  \item \textbf{KC frequency drift}: 2/17 KC có drift $> 3\%$ giữa train và test (KC 0: 3.9\%, KC 3: 3.5\%).
  \item \textbf{Class imbalance}: positive rate 82.5\% -- cần sử dụng PR-AUC thay vì chỉ ROC-AUC.
  \item \textbf{Quality distribution}: điểm self-rating phân phối đều trong $[0, 1]$.
  \item \textbf{Cold-start coverage}: 100\% người dùng có cold-start vector.
\end{enumerate}

\section{Cấu Hình Thực Nghiệm (Ablation Study)}
\label{sec:configs}

Để đánh giá tác động của từng yếu tố thiết kế, luận văn tiến hành ablation study với 8 cấu hình:

\begin{table}[H]
  \centering
  \caption{Các cấu hình trong ablation study}
  \label{tab:configs}
  \resizebox{\textwidth}{!}{%
  \begin{tabular}{clllp{5cm}}
    \toprule
    \textbf{\#} & \textbf{Tên} & \textbf{Model} & \textbf{Label} & \textbf{Siêu tham số} \\
    \midrule
    0 & Baseline    & Majority class & binary & -- \\
    1 & BKT         & Bayesian HMM   & binary & L-BFGS-B per KC \\
    2 & DKT-Binary  & LSTM           & binary & $h{=}128$, $L{=}1$, $d{=}0.2$, lr$=10^{-3}$ \\
    3 & DKT-Quality$\star$ & LSTM    & quality (MSE) & $h{=}128$, $L{=}1$, $d{=}0.2$, lr$=10^{-3}$ \\
    4 & SAKT-Binary & Self-Attention  & binary & $e{=}64$, $H{=}4$, $L{=}1$, $d{=}0.2$ \\
    5 & SAKT-Quality$\star$ & Self-Attention & quality (MSE) & $e{=}64$, $H{=}4$, $L{=}1$, $d{=}0.2$ \\
    6 & DKT-Deep    & LSTM (deep)    & binary & $h{=}256$, $L{=}2$, $d{=}0.3$, lr$=5\times10^{-4}$ \\
    7 & SAKT-Deep   & Self-Attention (deep) & binary & $e{=}128$, $H{=}4$, $L{=}2$, $d{=}0.2$, lr$=5\times10^{-4}$ \\
    \bottomrule
  \end{tabular}%
  }
\end{table}

Tất cả mô hình neural được huấn luyện với: batch size = 64, epochs = 30, early stopping patience = 5.

\section{Độ Đo Đánh Giá}
\label{sec:metrics}

\begin{itemize}
  \item \textbf{ROC-AUC}: đo khả năng phân biệt đúng/sai, không phụ thuộc ngưỡng.
  \item \textbf{PR-AUC} (Average Precision): phù hợp khi class imbalance cao (82.5\% positive).
  \item \textbf{PR-AUC Gain}: $\text{PR-AUC}_{\text{model}} - \text{PR-AUC}_{\text{baseline}}$, đo cải thiện thực sự so với đường cơ sở.
  \item \textbf{RMSE}: sai số bình phương trung bình, đánh giá độ chính xác dự đoán xác suất.
  \item \textbf{Accuracy}: tỷ lệ phân loại đúng với ngưỡng 0.5.
\end{itemize}

\section{Kết Quả Ablation Study}
\label{sec:results}

\begin{table}[H]
  \centering
  \caption{Kết quả ablation study trên tập test (150 người dùng, 6{,}065 tương tác)}
  \label{tab:results}
  \begin{tabular}{lrrrrrrr}
    \toprule
    \textbf{Model} & \textbf{ROC-AUC} & \textbf{PR-AUC} & \textbf{PR-AUC Gain} & \textbf{RMSE} & \textbf{Acc} & \textbf{Thời gian (s)} \\
    \midrule
    Baseline           & 0.5000 & 0.8201 & 0.0000 & 0.3841 & 0.8201 & 0.0 \\
    BKT                & \textbf{0.6767} & 0.9053 & 0.0852 & 0.3735 & 0.8201 & 54.8 \\
    DKT-Binary         & 0.6662 & 0.9048 & 0.0847 & 0.3751 & 0.8220 & 28.0 \\
    DKT-Quality$\star$ & 0.6720 & \textbf{0.9074} & \textbf{0.0873} & \textbf{0.3553} & 0.7956 & 19.4 \\
    SAKT-Binary        & 0.6684 & 0.9033 & 0.0832 & 0.3733 & 0.8220 & \textbf{11.0} \\
    SAKT-Quality$\star$& 0.6674 & 0.9039 & 0.0838 & 0.3564 & 0.7765 & 11.8 \\
    DKT-Deep           & 0.6693 & 0.9059 & 0.0858 & 0.3745 & 0.8220 & 110.0 \\
    SAKT-Deep          & 0.6695 & 0.9058 & 0.0857 & 0.3734 & 0.8220 & 30.3 \\
    \bottomrule
  \end{tabular}
\end{table}

\begin{figure}[H]
  \centering
  \includegraphics[width=\textwidth]{ablation_bar.png}
  \caption{So sánh ROC-AUC và PR-AUC của 8 cấu hình trên tập test}
  \label{fig:ablation_bar}
\end{figure}

\begin{figure}[H]
  \centering
  \includegraphics[width=\textwidth]{learning_curves.png}
  \caption{Learning curves của 6 mô hình neural (train loss, val loss, val AUC theo epoch)}
  \label{fig:learning_curves}
\end{figure}

\section{Phân Tích Kết Quả}
\label{sec:analysis}

\subsection{Hiệu quả chung của KT pipeline}

Tất cả 7 mô hình KT đều vượt Baseline một cách rõ rệt: ROC-AUC tăng từ 0.50 lên 0.67--0.68, PR-AUC Gain đạt 0.083--0.087. Điều này chứng minh rằng việc sử dụng lịch sử tương tác để dự đoán kết quả câu hỏi tiếp theo mang lại giá trị thực sự, ngay cả với dữ liệu mô phỏng.

\subsection{Ảnh hưởng của quality label}

DKT-Quality$\star$ đạt PR-AUC cao nhất (0.9074) và RMSE thấp nhất (0.3553), vượt DKT-Binary (PR-AUC 0.9048, RMSE 0.3751). Điều này xác nhận giả thuyết rằng quality label -- điểm tự đánh giá liên tục thay vì nhãn nhị phân -- cung cấp tín hiệu gradient phong phú hơn, giúp mô hình học được biểu diễn trạng thái kiến thức tốt hơn. Tuy nhiên, accuracy của DKT-Quality$\star$ thấp hơn (0.7956 so với 0.8220) do mô hình tối ưu hóa MSE thay vì cross-entropy, làm phân phối đầu ra lệch khỏi ngưỡng 0.5.

\subsection{BKT và simulation bias}

Đáng chú ý, BKT đạt ROC-AUC cao nhất (0.6767), vượt tất cả mô hình neural. Điều này được giải thích bởi \textbf{simulation bias}: dữ liệu được sinh theo mô hình IRT -- về cơ bản có cùng dạng tham số với BKT (per-KC parameterization). Nói cách khác, BKT và quá trình sinh dữ liệu chia sẻ cùng inductive bias, tạo lợi thế không công bằng cho BKT. Trên dữ liệu thực từ người dùng thật, kỳ vọng mô hình neural sẽ cho kết quả tốt hơn nhờ khả năng học được các pattern phức tạp hơn.

\subsection{Deep models và data size bottleneck}

DKT-Deep (2-layer LSTM, 256 hidden units) không cải thiện đáng kể so với DKT-Binary (ROC-AUC 0.6693 so với 0.6662), trong khi thời gian huấn luyện tăng gấp 4 lần (110s so với 28s). Tương tự, SAKT-Deep không vượt SAKT-Binary. Điều này cho thấy với 28{,}164 tương tác train, mô hình phức tạp hơn không mang lại lợi ích do \textbf{data size bottleneck} -- dữ liệu chưa đủ để khai thác năng lực của mô hình sâu hơn.

\subsection{Hiệu quả tính toán}

SAKT-Binary đạt thời gian huấn luyện ngắn nhất (11s) với ROC-AUC cạnh tranh (0.6684), nhờ khả năng song song hóa của self-attention. Đây là lựa chọn tốt cho môi trường production với yêu cầu update model thường xuyên.

\section{Phân Tích Sâu (Notebook 04)}
\label{sec:deep_analysis}

\subsection{Hiệu suất theo từng Knowledge Component}

\begin{figure}[H]
  \centering
  \includegraphics[width=\textwidth]{per_kc_auc.png}
  \caption{ROC-AUC theo từng KC: BKT vs DKT-Quality$\star$ trên tập test}
  \label{fig:per_kc_auc}
\end{figure}

Hình \ref{fig:per_kc_auc} cho thấy hiệu suất của mô hình không đồng đều giữa các KC. Các KC có tần suất thấp (KC~0: SQL\_DATABASE với drift 3.9\%, KC~3: với drift 3.5\%) có AUC thấp hơn do ít dữ liệu huấn luyện hơn. DKT-Quality$\star$ vượt BKT ở phần lớn các KC thuộc DS và DE, trong khi BKT cạnh tranh hơn ở các KC DA -- phù hợp với lý giải simulation bias (DA có skill phân bố đơn giản hơn, phù hợp với parametric form của BKT).

\subsection{Phân tích theo vị trí trong chuỗi}

\begin{figure}[H]
  \centering
  \includegraphics[width=0.85\textwidth]{position_auc.png}
  \caption{ROC-AUC theo vị trí trong sequence: early (0--25\%) vs late (75--100\%)}
  \label{fig:position_auc}
\end{figure}

Hình \ref{fig:position_auc} cho thấy cả ba mô hình (BKT, DKT-Quality$\star$, SAKT-Deep) đều có xu hướng cải thiện AUC khi có nhiều lịch sử hơn (early $\to$ late). Điều này xác nhận rằng mô hình KT học tốt hơn khi có đủ ngữ cảnh -- một lý do quan trọng để thu thập nhiều dữ liệu tương tác từ mỗi người dùng.

\subsection{Phân tích theo độ khó câu hỏi}

\begin{figure}[H]
  \centering
  \includegraphics[width=0.85\textwidth]{difficulty_auc.png}
  \caption{ROC-AUC phân tầng theo độ khó: EASY / MEDIUM / HARD}
  \label{fig:difficulty_auc}
\end{figure}

Như thể hiện trong Hình \ref{fig:difficulty_auc}, câu hỏi MEDIUM có AUC cao nhất -- đây là vùng có độ phân biệt cao nhất trong IRT. Câu hỏi EASY bị ảnh hưởng bởi class imbalance nặng hơn (hầu hết trả lời đúng), trong khi HARD có ít dữ liệu hơn.

\subsection{Phân tích theo loại người dùng}

\begin{figure}[H]
  \centering
  \includegraphics[width=0.85\textwidth]{usertype_auc.png}
  \caption{ROC-AUC theo loại người dùng: beginner / intermediate / advanced}
  \label{fig:usertype_auc}
\end{figure}

Hình \ref{fig:usertype_auc} cho thấy người dùng intermediate có AUC cao nhất, trong khi beginner và advanced thấp hơn. Điều này phù hợp với lý thuyết IRT: người dùng ở mức trung bình có phân phối đúng/sai cân bằng hơn, cung cấp tín hiệu học tốt hơn cho mô hình.

\subsection{Cold-start score và độ chính xác mô hình}

\begin{figure}[H]
  \centering
  \includegraphics[width=0.7\textwidth]{coldstart_vs_acc.png}
  \caption{Tương quan giữa cold-start score và per-user AUC (DKT-Quality$\star$)}
  \label{fig:coldstart_vs_acc}
\end{figure}

Hình \ref{fig:coldstart_vs_acc} kiểm tra giả thuyết rằng người dùng có cold-start vector chất lượng cao (self-rating tốt) sẽ được dự đoán chính xác hơn. Hệ số tương quan Pearson $r$ cho thấy mối tương quan dương yếu, xác nhận cold-start score mang thông tin bổ sung nhưng không phải yếu tố duy nhất ảnh hưởng đến độ chính xác mô hình.

\subsection{Phân tích lỗi}

\begin{figure}[H]
  \centering
  \includegraphics[width=\textwidth]{error_analysis.png}
  \caption{Confusion matrix và phân phối điểm dự đoán theo nhãn thực -- DKT-Quality$\star$}
  \label{fig:error_analysis}
\end{figure}

Hình \ref{fig:error_analysis} (trái) cho thấy mô hình có xu hướng dự đoán Pass (positive) do class imbalance -- False Negative Rate thấp hơn False Positive Rate. Phân phối điểm dự đoán (phải) cho thấy mô hình phân biệt được hai nhóm, nhưng có sự chồng lấp đáng kể do nhiễu trong dữ liệu mô phỏng.

\section{Hạn Chế và Thảo Luận}
\label{sec:limitations}

\subsection{Simulation bias (nghiêm trọng)}
Dữ liệu sinh theo IRT tạo lợi thế cho BKT vì cùng inductive bias. Kết quả thực nghiệm trên dữ liệu người dùng thật có thể khác biệt đáng kể -- đây là hạn chế quan trọng nhất cần được giải quyết trong nghiên cứu tương lai bằng cách thu thập dữ liệu thực từ người dùng InternHub.

\subsection{Class imbalance (trung bình)}
Positive rate 82.5\% làm PR-AUC bị inflate. PR-AUC Gain ($\approx 0.085$) là độ đo đáng tin cậy hơn để so sánh giữa các mô hình.

\subsection{KC frequency imbalance (nhẹ)}
2/17 KC có drift $> 3\%$ giữa train/test, ảnh hưởng nhỏ đến tính tổng quát. Có thể giải quyết bằng stratified split theo KC trong tương lai.

\subsection{Data scale (nhẹ)}
28K tương tác train là đủ cho BKT và shallow neural models, nhưng chưa đủ để khai thác deep models (2-layer LSTM/Transformer). Cần ít nhất 100K tương tác để thấy lợi ích rõ ràng từ kiến trúc sâu hơn.

\subsection{Sequence length hằng định (nhẹ)}
Mặc dù pipeline thiết kế randomize số câu hỏi mỗi người dùng trong khoảng 20--60, trên thực tế toàn bộ 1{,}000 người dùng đều có đúng 40 tương tác. Nguyên nhân là hàm \texttt{stratified\_sample\_questions} trả về toàn bộ pool khi \texttt{len(pool) <= n} -- pool câu hỏi mỗi role sau bộ lọc competency chỉ có $\sim$40 câu duy nhất, khiến tham số randomization không có tác dụng. Hậu quả là mô hình KT không được huấn luyện trên các sequence có độ dài khác nhau, làm giảm khả năng tổng quát hóa sang người dùng thực (có session length thay đổi nhiều hơn). Giải pháp trong tương lai: mở rộng question bank lên $\geq$100 câu/role.

\subsection{Phân phối trình độ người dùng lệch (nhẹ)}
Phân phối thực tế sau reclassification là beginner = 40.9\%, intermediate = 53.2\%, advanced = 5.9\%, lệch so với mục tiêu thiết kế (35\%/50\%/15\%). Nhóm advanced bị thiếu hụt đáng kể do base\_skill = 0.72 của tier advanced, khi cộng noise $\mathcal{N}(0, 0.15)$, thường cho mean skill nằm trong khoảng $[0.40, 0.70]$ và bị phân loại lại thành intermediate. Hậu quả là mô hình ít được huấn luyện trên người dùng có skill cao, có thể ảnh hưởng đến hiệu suất dự đoán cho nhóm advanced trong thực tế.

\section{Tóm Tắt Kết Quả}
\label{sec:summary}

Từ các thực nghiệm, rút ra các kết luận chính:

\begin{enumerate}
  \item \textbf{KT pipeline có giá trị}: tất cả mô hình KT vượt Baseline (ROC-AUC 0.67 so với 0.50), PR-AUC Gain $\approx 0.085$.
  \item \textbf{Quality label cải thiện DKT}: DKT-Quality$\star$ đạt PR-AUC và RMSE tốt nhất trong nhóm neural models.
  \item \textbf{BKT vượt trội do simulation bias}: không phản ánh hiệu năng thực tế trên dữ liệu người dùng thật.
  \item \textbf{Deep models chưa hiệu quả}: bottleneck là data size, không phải kiến trúc.
  \item \textbf{SAKT-Binary hiệu quả nhất về tính toán}: AUC cạnh tranh, thời gian train chỉ 11s.
  \item \textbf{Mô hình học tốt hơn với nhiều lịch sử}: AUC tăng dần theo vị trí trong sequence.
\end{enumerate}

\textbf{Khuyến nghị cho production}: sử dụng DKT-Quality$\star$ làm mô hình chính (cân bằng giữa accuracy và interpretability), kết hợp SAKT-Binary cho online update nhanh.
